\chapter*{Annexes}
\addcontentsline{toc}{chapter}{Annexes}

\section*{Annexe 1 : Extrait du Dictionnaire de Données (MCD)}
\begin{table}[H]
    \centering
    \begin{tabular}{|l|l|l|l|}
        \hline
        \textbf{Nom de la donnée} & \textbf{Code} & \textbf{Type} & \textbf{Taille} \\
        \hline
        Identifiant Véhicule & \texttt{VEH\_ID} & Numérique & Auto-incrément \\
        Immatriculation & \texttt{VEH\_IMMAT} & Alphanumérique & 15 \\
        Prix de location & \texttt{VEH\_PRIX} & Décimal & - \\
        Date de début de réservation & \texttt{RESA\_DATE\_DEBUT} & Date & - \\
        \hline
    \end{tabular}
    \caption{Extrait du dictionnaire de données}
\end{table}

\section*{Annexe 2 : Extrait de la Documentation de l'API REST}
Au lieu d'inclure le code brut, voici un extrait de la documentation de l'API (de type Swagger/OpenAPI) exposée par les microservices backend à destination du portail React :
\begin{table}[H]
    \centering
    \renewcommand{\arraystretch}{1.3}
    \begin{tabular}{|l|l|p{8cm}|}
        \hline
        \textbf{Méthode HTTP} & \textbf{Endpoint (URL)} & \textbf{Description / Action métier} \\
        \hline
        \texttt{GET} & \texttt{/api/v1/vehicles} & Récupère la liste des véhicules (avec cache Redis) \\
        \texttt{GET} & \texttt{/api/v1/vehicles/:id} & Détails d'un véhicule spécifique \\
        \texttt{POST} & \texttt{/api/v1/auth/login} & Authentification et génération du Token JWT \\
        \texttt{POST} & \texttt{/api/v1/bookings} & Création d'une réservation (Stripe Webhook) \\
        \texttt{GET} & \texttt{/api/v1/erp/invoices} & Génération des factures PDF (RBAC: Admin) \\
        \hline
    \end{tabular}
    \caption{Extrait des routes de l'API RESTful}
\end{table}

\section*{Annexe 3 : Charte Graphique et Expérience Utilisateur (UI/UX)}
L'interface de la plateforme B2C a été conçue pour refléter le positionnement haut de gamme de \textbf{Prestige Maroc}. Voici les principes de la charte graphique :
\begin{itemize}
    \item \textbf{Typographie :} Utilisation de polices modernes et sans-serif (ex: \textit{Inter, Roboto}) pour maximiser la lisibilité sur les appareils mobiles.
    \item \textbf{Palette de Couleurs :}
    \begin{itemize}
        \item \textbf{Couleur Primaire :} Noir profond (\#121212) évoquant le luxe et l'élégance.
        \item \textbf{Couleur Secondaire :} Or / Doré (\#D4AF37) pour les boutons d'appels à l'action (Call-To-Action).
        \item \textbf{Couleur d'Arrière-plan :} Blanc pur (\#FFFFFF) ou Gris clair (\#F8F9FA) pour faire ressortir les photographies des véhicules.
    \end{itemize}
    \item \textbf{Composants interactifs :} Effets de "Glassmorphism" (verre dépoli) et ombres portées (Drop Shadows) générés via Tailwind CSS.
\end{itemize}

\section*{Annexe 4 : Diagramme de Gantt (Export MS Project)}
\textit{(Le jury trouvera ici l'impression au format A3 du diagramme de Gantt complet du projet, s'étalant du mois de Mars au mois de Juin 2026).}
